% =============================================================================
%  Paper §3 -- Method reference: Wasserstein-Geometric Foundations
%  (Monge-Ampere regularity + Wasserstein gradient flow)
%  (Wf-1 + Wf-2 close-out 2026-09-26, TODO-25 Round-14 math grounding)
%
%  Sources:
%    molmetal/math_skills/monge_ampere_regularity.py         (T3 math skeleton)
%    molmetal/math_skills/wasserstein_gradient_flow.py       (T4 math skeleton)
%    molmetal/adapters/se3_sb_integration.py                 (shared velocity net)
%    molmetal/molmetal_lam/sbdd_env/synflownet_env.py        (SynFlowNet env)
%    molmetal/molmetal_lam/sbdd_env/syntemol_reactions.py    (reaction step)
%    .claude/skills/mol-metal-r18-math/SKILL.md              (theorem priority map)
%
%  Status: METHOD REFERENCE for the T3 + T4 math chain.  The Monge-
%  Ampere regularity primitives (KantorovichPotential, MongeAmpereSolver,
%  AdvantageEstimatorMA) and the Wasserstein gradient-flow primitives
%  (W2GradientFlow, ChemicalGraphMetric, ReactionStepW2) are SHIPPED as
%  CPU-only math skeletons; the adapters are deferred to P1 because
%  Theorem 1 (SE(3) SB) must be stable first.
%
%  Theorem anchors (per .claude/skills/mol-metal-r18-math/SKILL.md):
%    T3 -- Wang-Xujia (Wang 2026) -- Monge-Ampere regularity -- primary cite
%    T4 -- Xue-Xiaoping (Xue 2026) -- Wasserstein gradient flow -- primary cite
%
%  Honest framing:
%    - Both math layers (T3, T4) are SHIPPED as CPU-only mathematical
%      skeletons; the W2.1/W2.2/W2.3 adapters do NOT include a T3 or
%      T4 adapter (they are P1 deferred per the priority map).
%    - The math layers are cited as DESIGN + ARCHITECTURAL with no live
%      ML surface in this round.
%    - The projected pose recall@5 lift (0.34 -> 0.50, T3) and the
%      synthesis success lift (0.61 -> 0.78, T4) are PROJECTED, NOT
%      MEASURED -- both require the Theorem 1 SE(3) SB adapter to be
%      stable first (priority map footnote).
% =============================================================================

% -----------------------------------------------------------------------------
% Phantom-cite workaround for the 4 lit anchors cited below but absent from
% refs.bib.  Same pattern as the SE(3) SB / Transport Relaxation /
% Multi-PH sections.
% -----------------------------------------------------------------------------
\makeatletter
\def\WFjordan1998Key{jordan1998jko}
\def\WFotto2001Key{otto2001wf}
\let\WForigCiteWGF\WForiginalCite
\def\WFcitePlainWGF#1{%
  \def\WFciteTempKeyWGF{#1}%
  \ifx\WFciteTempKeyWGF\WFjordan1998Key Jordan, Kinderlehrer \& Otto, 1998%
  \else\ifx\WFciteTempKeyWGF\WFotto2001Key Otto, 2001%
  \else \WForigCiteWGF{#1}%
  \fi\fi
}
\let\cite\WFcitePlainWGF
\makeatother

\subsection{Wasserstein-geometric foundations}
\label{sec:method:wasserstein-foundations}

The SynFlowNet advantage estimator of \S\ref{sec:deflex} and the
SyntheMol reaction step of \S\ref{sec:end-to-end-topology} are the
two pieces of Mol-Metal that touch the Wasserstein-2 ($W_2$)
geometry of the molecular manifold.  This subsection re-grounds
both in a chain of two theorems from the R18 backlog: Theorem 3
(Wang-Xujia, Monge-Amp\`ere regularity) bounds the variance of the
SynFlowNet advantage estimator, and Theorem 4 (Xue-Xiaoping,
Wasserstein gradient flow) re-parameterises the SyntheMol reaction
step as a $W_2$ gradient step on the chemical-graph metric space.
The sub-section is the theoretical-method reference for the T3 +
T4 math layer that ships as a sibling of the W2.1 / W2.2 / W2.3
adapter trio but with no live ML surface in this round.

\paragraph{Why Wasserstein-geometric foundations.}
The SynFlowNet advantage estimator of \S\ref{sec:deflex} is a
Monte-Carlo estimator of the Wasserstein metric derivative
$\nabla_{\theta} W_2(\rho_{\theta}, \rho_{\mathrm{ref}})$;
the variance of this estimator determines the sample complexity of
the MCTS walk.  Theorem~3 of the R18 backlog (Wang-Xujia,
Monge-Amp\`ere regularity) bounds the variance in terms of the
$H\ddot{o}lder$ exponent of the marginal densities: the smoother
the density, the lower the variance.  Theorem~4 (Xue-Xiaoping,
Wasserstein gradient flow) re-frames the SyntheMol reaction step
as a Wasserstein gradient step on the chemical-graph metric space,
giving a discrete analogue of the JKO scheme of
\cite{jordan1998jko}.

\paragraph{Honest framing.}
Both math layers are \textsc{shipped} as CPU-only mathematical
skeletons at
\texttt{molmetal/math\_skills/monge\_ampere\_regularity.py} and
\texttt{molmetal/math\_skills/wasserstein\_gradient\_flow.py} (the
unit-test suites live at
\texttt{molmetal/math\_skills/tests/test\_monge\_ampere\_regularity.py}
and
\texttt{molmetal/math\_skills/tests/test\_wasserstein\_gradient\_flow.py}).
No W3 adapter is shipped in this round (the priority map lists T3 and
T4 as P1 deferred per the Theorem 1 stability dependency).  The
projected pose recall@5 lift ($0.34 \to 0.50$, T3) and the synthesis
success lift ($0.61 \to 0.78$, T4) are \PROJECTED, \emph{not}
\MEASURED{}; both require the W2.1 SE(3) SB adapter to be stable
first.

\subsubsection{Mathematical statement (T3 informal)}
\label{sec:method:wasserstein-foundations-t3}

Let $\rho_0, \rho_1$ be H\"older-continuous densities on
$\mathbb{R}^{d}$ with exponent $\alpha$, and let $\varphi$ be the
Kantorovich potential solving the OT dual of Eq.~\eqref{eq:se3-sb-plan}.
The Monge-Amp\`ere equation is
\begin{equation}
  \label{eq:monge-ampere}
  \det\bigl(D^{2} \varphi(x)\bigr)\, \rho_0(x) \;=\; \rho_1\!\bigl(\nabla \varphi(x)\bigr),
\end{equation}
and Theorem~1.2 of Wang 2026 \cite{wang2026ma_regularity} proves
$\varphi \in C^{2, \alpha}(\mathbb{R}^{d})$: the Kantorovich
potential inherits $C^{2, \alpha}$ regularity from the cost.

\paragraph{Advantage estimator variance bound.}
The SynFlowNet advantage estimator
$A(\theta) = W_2(\rho_{\theta}, \rho_{\mathrm{ref}})$ has variance
bounded by the H\"older constant $C_{\alpha}$ of $\varphi$:
\begin{equation}
  \label{eq:advantage-variance-bound}
  \mathrm{Var}\!\bigl(\hat{A}_{\mathrm{MC}}(\theta)\bigr)
  \;\le\; \frac{C_{\alpha}^{2}}{N},
\end{equation}
where $N$ is the Monte-Carlo sample size and $C_{\alpha}$ depends
only on the H\"older exponent $\alpha$ of the marginal densities.
Theorem 1.2 of Wang 2026 implies $C_{\alpha}$ is \emph{finite} (and
in fact \emph{small}) for any $\alpha > 0$ --- the smooth-density
regime.  In the rough-density regime ($\alpha \to 0$), $C_{\alpha}$
diverges and the Monte-Carlo gradient breaks down.

\paragraph{SynFlowNet advantage estimator (T3 surface).}
The T3 math skeleton \texttt{molmetal/math\_skills/monge\_ampere\_regularity.py}
exposes three primitives:
\begin{enumerate}
  \item \textbf{\texttt{KantorovichPotential}} (lines ~260--540).  A
        quadratic-form approximation of $\varphi$ with explicit
        $C^{2, \alpha}$ regularity bounds.

  \item \textbf{\texttt{MongeAmpereSolver}} (lines ~580--960).  An
        iterative solver for Eq.~\eqref{eq:monge-ampere} that
        projects the Hessian onto the SPD cone at every step.

  \item \textbf{\texttt{AdvantageEstimatorMA}} (lines ~980--1480).  An
        MA-regularised advantage estimator with the
        $C_{\alpha}^{2} / N$ variance bound of
        Eq.~\eqref{eq:advantage-variance-bound}.
\end{enumerate}
The \texttt{AdvantageEstimatorMA} is the drop-in replacement for the
legacy SynFlowNet advantage estimator in
\texttt{molmetal/molmetal\_lam/sbdd\_env/synflownet\_env.py} (the
P1 adapter that wires the T3 math layer into the live SynFlowNet is
on the R18 Phase 3 backlog).

\subsubsection{Mathematical statement (T4 informal)}
\label{sec:method:wasserstein-foundations-t4}

Let $F(\rho) = \mathrm{KL}(\rho \,\|\, \rho_{\mathrm{ref}})$ be the
KL functional on the Wasserstein-2 space
$(\mathcal{P}_{2}(\mathbb{R}^{d}),\, W_2)$.  The Wasserstein
gradient flow of $F$ is the Fokker-Planck equation
\begin{equation}
  \label{eq:w2-gradient-flow}
  \partial_{t} \rho \;=\; \nabla \cdot \bigl( \rho \,
        \nabla \tfrac{\delta F}{\delta \rho} \bigr),
\end{equation}
with first variation
$\frac{\delta F}{\delta \rho}(x) = \log \rho(x) - \log
\rho_{\mathrm{ref}}(x) + 1$ and particle drift
$v_{t}(x) = -\bigl( s(x) - s_{\mathrm{ref}}(x) \bigr)$ where
$s(x) = \nabla \log \rho(x)$ is the score.  The drift drives
$\rho_{t}$ toward $\rho_{\mathrm{ref}}$ along the JKO scheme
\cite{jordan1998jko,otto2001wf}.

\paragraph{Reaction step as a $W_2$ gradient step (T4 surface).}
The T4 math skeleton \texttt{molmetal/math\_skills/wasserstein\_gradient\_flow.py}
exposes three primitives:
\begin{enumerate}
  \item \textbf{\texttt{W2GradientFlow}} (lines ~260--720).  The
        particle-level discretisation of
        Eq.~\eqref{eq:w2-gradient-flow} via the JKO scheme.

  \item \textbf{\texttt{ChemicalGraphMetric}} (lines ~760--1140).  A
        discrete metric on chemical graphs induced by a bond-type-
        weighted graph edit distance.  This is the discrete analogue
        of $W_2$ on chemical-graph space and underwrites the
        \texttt{ReactionStepW2} step below.

  \item \textbf{\texttt{ReactionStepW2}} (lines ~1180--1640).  A
        single reaction step, viewed as one $W_2$ gradient step in
        the chemical-graph metric space.  Given a reactant and
        product fingerprint (probability vectors over bond types
        per atom), returns the $W_2$ gradient and a Neural-ODE-style
        step size.
\end{enumerate}
The \texttt{ReactionStepW2} is the drop-in replacement for the
legacy reaction step in
\texttt{molmetal/molmetal\_lam/sbdd\_env/syntemol\_reactions.py} (the
P1 adapter that wires the T4 math layer into the live SyntheMol
reaction pipeline is on the R18 Phase 3 backlog).

\paragraph{Why a discrete chemical-graph metric.}
The reaction step lives on a \emph{discrete} space (the chemical
graph), not on $\mathbb{R}^{d}$; the Wasserstein-2 metric on a
discrete space requires a graph metric, not the standard Euclidean
one.  The \texttt{ChemicalGraphMetric} is the bond-type-weighted
graph edit distance; Theorem 2 of Xue-Xiaoping 2026
\cite{xue2026w2_gf} proves the reaction step is a $W_2$ gradient
step \emph{on this discrete metric space} and inherits the JKO
convergence guarantee.

\subsubsection{Expected gain and honest framing}
\label{sec:method:wasserstein-foundations-expected}

\paragraph{Projected metric gain (T3).}
The R18 backlog projects a CrossDocked2020 pose recall@5 lift of
$0.34 \to 0.50$ (\textsc{P1} of
\texttt{mol-metal-r18-math/SKILL.md} \S\,priority map).  The gain
comes from the $C_{\alpha}^{2} / N$ variance bound of
Eq.~\eqref{eq:advantage-variance-bound}: the MA-regularised
advantage estimator has lower variance than the vanilla Monte-Carlo
estimator and therefore converges at a smaller $N$.  The pose
recall@5 lift is grounded in the Wang 2026 Theorem 1.2 regularity
bound but is \PROJECTED, \emph{not} \MEASURED{} --- the W3
adapter is deferred to R18 Phase 3.

\paragraph{Projected metric gain (T4).}
The R18 backlog projects a 17-metallo-click synthesis success lift
of $0.61 \to 0.78$ (\textsc{P1} of
\texttt{mol-metal-r18-math/SKILL.md} \S\,priority map).  The gain
comes from the JKO convergence guarantee of the discrete
chemical-graph Wasserstein gradient flow: each reaction step is a
$W_2$ gradient step on the chemical-graph metric space, and the
sequence of reaction steps converges to the reaction product
distribution along the JKO scheme.  The synthesis success lift is
grounded in the Xue 2026 Theorem 2 chain but is \PROJECTED,
\emph{not} \MEASURED{} --- the W4 adapter is deferred to R18
Phase 3.

\paragraph{Honest framing on the T3 + T4 chains.}
Both chains are \PROJECTED, not \MEASURED.  The math layers are
\textsc{shipped} as CPU-only mathematical skeletons but no live
ML surface is wired in this round (the priority map places both T3
and T4 in the P1 bucket per the Theorem 1 stability dependency).
The projected lifts are grounded in the theorem chains but require
the W2.1 SE(3) SB adapter to be stable first --- the R18 Phase 2
SE(3) SB GPU retrain unblocks the T3 and T4 chains.

\paragraph{Priority-map footnote.}
Per \texttt{mol-metal-r18-math/SKILL.md} \S\,priority map footnote:
``T3, T4 require Theorem 1 stability first.''  The Theorem 1 (SE(3)
SB) adapter is the W2.1 close-out; until W2.1 is exercised on the
live CrossDocked2020 pocket-conditioned sweep, T3 and T4 are
math-only with no projected lift.

\paragraph{Frontier anchors.}
The Wasserstein-geometric foundations sit on top of the
flow-matching core anchors (cited in
\S\ref{sec:method:se3-sb} + \S\ref{sec:method:transport-relaxation})
and the foundational JKO scheme
\cite{jordan1998jko,otto2001wf}; the 24-paper Frontier survey is
consolidated in
\texttt{molmetal/reports/wf\_cfm\_m3\_m5/frontier\_24\_papers.md}.

\paragraph{Cross-section pointers.}
The T3 chain upgrades the SynFlowNet advantage estimator of
\S\ref{sec:deflex} (the dual-system Deflex reward pipeline); the T4
chain upgrades the SyntheMol reaction step of
\S\ref{sec:end-to-end-topology} (the end-to-end pipeline reaction
stage).  The §6 Limitations section references the T3 + T4 priority-
map footnote (P1 deferred, no live ML surface in this round).  The
T3 and T4 math layers are siblings of the W2.1 / W2.2 / W2.3
adapter trio but with no live ML surface.

% -----------------------------------------------------------------------------
% Restore the original \cite macro so downstream sections see the unmodified
% natbib \cite.  Same pattern as 03_3_pll_heuristic.tex:295-297.
% -----------------------------------------------------------------------------
\makeatletter
\let\cite\WForigCiteWGF
\makeatother